% Lösungen Einheit 3 von 3 – Wurzelgesetze und rationale Exponenten (zusammenbau v0.1)
\einheitenkopf{Einheit 3 von 3 · Wurzelgesetze und rationale Exponenten}

%% TODO Lösung der Erklärzeile 17g) fehlt in der Bank
\erg{17}{a) Produkt oder Quotient: trennbar \quad b) $\sqrt{576} = 24$ und $3 \cdot 8 = 24$ – gleich \quad c) $\sqrt{484} = 22$ und $2 \cdot 11 = 22$ – gleich \quad d) $\sqrt{900} = 30$ und $3 \cdot 10 = 30$ – gleich \quad e) $\sqrt{5\,184} = 72$ und $8 \cdot 9 = 72$ – gleich \quad f) $\sqrt{1\,089} = 33$ und $11 \cdot 3 = 33$ – gleich \quad g) \ldots \quad h) $\frac{8}{3}$ \quad i) $\sqrt{4 \cdot 7} = 2\sqrt{7} \approx 5{,}29$ \quad j) $6\sqrt{3} \approx 10{,}39$ \quad k) $\sqrt{225} = 15$ \quad l) $\sqrt{81} = 9$ \quad m) $4$ \quad n) $5^{\frac{1}{2}+\frac{1}{2}} = 5^1 = 5$ \quad o) $\frac{\sqrt{3}}{3} \approx 0{,}58$ \quad p) $a^{\frac{1}{2}} \cdot b^{\frac{1}{2}} = (a \cdot b)^{\frac{1}{2}}$ – gleicher Exponent, die Basen werden malgenommen. \quad q) $x = \sqrt[3]{512} = 8$ \quad r) $3\sqrt{2} + 2\sqrt{2} = 5\sqrt{2} = 5 \cdot 2^{\frac{1}{2}}$}
\erg{18}{a) Wurzeln aus Summanden darf man nicht zu einer Wurzel zusammenziehen: $\sqrt{81} + \sqrt{144} = 9 + 12 = 21$. \quad b) Etwa $a = 81$, $b = 144$: $\sqrt{225} = 15$, aber $9 + 12 = 21$. \quad c) $x^{\frac{3}{5}}$ \quad d) $a = \sqrt{400 \cdot 2}\,\text{m} = 20\sqrt{2}\,\text{m} \approx 28{,}3\,\text{m}$}
